\manualchapter{Troubleshooting and further reading}{sec:troubleshooting}
\begin{description}
\item[Nothing changes in the game.]
Put Genie immediately to the right of RackNES, load a game, select the
matching map, and assign the connected row. Check that the scene uses the
chosen variable. Some values are quickly overwritten by game logic.
\item[An Unassigned row or menu causes a crash.]
Select a game before opening any row menu or using Randomize. Keep
Unassigned rows unplugged. The current implementation has unchecked
selection access; these are required operating precautions for this version.
\item[Changing games produces strange labels or behavior.]
Map changes retain the previous indices. Unplug inputs and clear selections
before switching maps, then choose entries from the new map.
\item[A toggle appears to ignore the gate.]
Use distinct rising edges above about 2 V, separated by lows below about
0.1 V. A high gate writes only on its rising edge. The row's internal state
may disagree with the game, and the game can overwrite the byte.
\item[Only 0 and 1 are written despite a larger listed range.]
Entries marked T use the toggle path regardless of their endpoint values.
Spawn Powerup, Reverb (Pulse), and Potion Type are affected.
\item[LOAD does not undo a change.]
Disconnect continuous CV writes before restoring RackNES's snapshot. Genie
selections and toggle states are outside that snapshot.
\item[Two controls change together.]
Check for duplicate addresses. The Mario Enemy Heading entries all share
\texttt{0x001A}, also used by Enemy 5 Type. Row order matters for competing
writes.
\item[Values wrap or behave unpredictably.]
Check that continuous CV stays in 0--10 V. There is no input clamp. Also
check whether the selected byte value is meaningful in the current game
scene; the table endpoints are not validation of every possible state.
\end{description}

\subsection{Sources and credits}
CV Genie was developed by \textbf{@anlexmatos}. The embedded maps credit
Data Crystal. The tables in this manual reproduce the project's
\texttt{src/GameMaps.hpp}; operating behavior follows \texttt{src/CVGenie.cpp}
and the expander handling in \texttt{src/RackNES.cpp}. They describe this
implementation rather than corrected or expanded external memory maps.
See the \href{https://github.com/Kautenja/RackNES/releases/latest/download/RackNES.pdf}{RackNES manual}
for ROM loading, controller gates, snapshots, clocking, and audio routing.
